% Sección incorporable; requiere amsmath y amssymb.
% Fuente principal: RIGIDEZ_SIMULTANEA_DEL_LECTOR_RADIAL_20260926.md.
% Complementos demostrativos: COMPOSICION_METRICA_MEMORIA_Y_ACOPLAMIENTO.md
% y ACCION_ACOPLADA_FASE_Y_MEMORIA_DE_FRONTERA.md.
% La integración, la compilación y la revisión visual corresponden al editor.
\section{Rigidez simultánea del lector longitudinal y del acoplamiento radial}
\label{g26:sec:rigidez-radial-es}

La construcción de Holografía Modular Triádica transmite a la realización
radial un estado enriquecido con sus coordenadas, orientación, acarreo,
memoria e incidencias. La sucesión
\[
 \begin{gathered}
 \mathrm{APP}\longrightarrow\mathrm{TRIT}\longrightarrow\mathrm{TPK}
 \longrightarrow\text{estado enriquecido}\\
 \longrightarrow\text{estructura discreta conjunta del continuo}
 \end{gathered}
\]
conserva la proyección arquimediana y la proyección de incidencia, junto con
las cinco construcciones consustanciales del continuo. La prolongación
\[
 w_6\longrightarrow w_{12}\longrightarrow w_{18}\longrightarrow
 w_{24}\longrightarrow w_{30}\longrightarrow R_{36}\longrightarrow G_9
\]
se articula con el orden operatorio
\(U_t\longrightarrow\Gamma_9\longrightarrow K_{\rm ph}\): la fase retorna
y el registro de memoria se incrementa. Las coordenadas regionales, el
registro seleccionado \(K\), la coordenada \(\alpha\) y la acción de retorno
se reciben con esta procedencia común y con sus lectores previamente
construidos.

El resultado de esta sección es una clasificación de realizaciones que
conservan simultáneamente los datos generados y las reglas longitudinales,
métricas y temporales que se precisan a continuación. Para cada carácter
positivo queda determinada la respuesta radial; el coeficiente que relaciona
esa respuesta con la masa es común a todos los caracteres admisibles.
La interpretación gravitatoria se formula después de construir ambas
lecturas sobre el mismo espacio.

\subsection{Registro de incidencias y restricciones constitutivas}

Se consideran espacios de Hilbert complejos \(V_{\rm in}\) y
\(V_{\rm out}\) de dimensión \(N\), con sus productos escalares fijados.
El cardinal marcado que se especifica en R1 es positivo; en particular,
ambas fibras son no nulas. Los adjuntos se toman respecto de estos productos. El símbolo
\(\mathsf U\) designará el transporte normalizado del archivo y
\(U_t\) conservará su significado de evolución del estado enriquecido.

\paragraph{R1. Registro marcado.}
El conjunto de incidencias es
\begin{equation}
 \Omega=\left\{(j,k,q,u,v):
 \begin{array}{l}
 j,k\in\mathbb Z/12\mathbb Z,\quad k-j\in\{0,3,6,9\},\\
 1\leq q\leq8,\quad u<v,\quad
 u,v\in\{1,2,4,5,7,8\}
 \end{array}\right\}.
 \label{g26:eq:omega-rigidez-es}
\end{equation}
Las doce fases, las cuatro posiciones relativas, las ocho posiciones del
índice \(q\) y las quince parejas del último conjunto dan
\begin{equation}
 N=|\Omega|=12\cdot4\cdot8\binom62
   =12\cdot4\cdot120=5760.
 \label{g26:eq:cardinal-rigidez-es}
\end{equation}
Cada incidencia posee su marca individual. En la base correspondiente
\((e_i)_{i\in\Omega}\) de \(V_{\rm out}\), sean
\(P_i=e_ie_i^*\). El modo común \(\boldsymbol u\) está normalizado por
\(|u_i|^2=1/N\); se escribe \(P=\boldsymbol u\boldsymbol u^*\).
El álgebra marcada es
\(\mathcal A_\Omega=\operatorname{span}\{P_i:i\in\Omega\}\).

\paragraph{R2. Normalización y archivo completo.}
El inversor generado \(B:V_{\rm in}\to V_{\rm out}\) satisface
\begin{equation}
 B^*B=\tfrac14 I_{\rm in},\qquad
 BB^*=\tfrac14 I_{\rm out}.
 \label{g26:eq:normalizacion-b-es}
\end{equation}
Se conservan las dos componentes
\begin{equation}
 C_0=(I-P)B,\qquad M_0=PB,\qquad
 C_0+M_0=B,\qquad \mathsf U=2B.
 \label{g26:eq:archivo-completo-es}
\end{equation}
La primera registra la variación respecto del modo común; la segunda
conserva dicho modo. Sus imágenes son ortogonales y su suma recupera
el transporte completo. Por \eqref{g26:eq:normalizacion-b-es},
\(\mathsf U\) es unitario.

\paragraph{R3. Inscripción individual.}
La aplicación
\(\mathcal E:\operatorname{End}(V_{\rm out})\to\mathcal A_\Omega\)
es una expectativa condicional lineal que fija cada \(P_i\) y satisface
\begin{equation}
 \mathcal E(A_1TA_2)=A_1\mathcal E(T)A_2,
 \qquad A_1,A_2\in\mathcal A_\Omega.
 \label{g26:eq:bimodularidad-es}
\end{equation}
La inscripción del resultado conserva, en un registro complementario,
los datos completos de la construcción. La expectativa determina el
observable inscrito; la conservación del archivo determina su
reconstrucción posterior.

\paragraph{R4. Composición longitudinal.}
Se fijan la coordenada generada \(\alpha>0\) y la sección de referencia
longitudinal \(L_*>0\). El transporte longitudinal se define por
\begin{equation}
 D=L_*\alpha^{16}\mathcal E(C_0C_0^*)\mathsf U:
 V_{\rm in}\longrightarrow V_{\rm out}.
 \label{g26:eq:regla-longitudinal-es}
\end{equation}
Esta regla aplica linealmente el coeficiente inscrito a la sección
longitudinal. Su tipo fija la posición de la potencia \(\alpha^{16}\)
y del coeficiente de retorno dentro de la composición. La lectura de
una raíz de intensidad constituye una operación diferente y posee su
propia regla de transformación.

\paragraph{R5. Métrica radial positiva.}
Sea \(X=X^*>0\) el carácter completo sobre \(V_{\rm in}\). Se define
la lectura radial positiva \(R_X\) sobre \(V_{\rm out}\) mediante
\begin{equation}
 \|R_Xv\|^2=\|XD^*v\|^2
 \qquad(v\in V_{\rm out}).
 \label{g26:eq:metrica-radial-es}
\end{equation}
Esta igualdad fija la métrica en cada vector de la fibra. En notación
polar, \(R_X\) es el módulo de llegada
\(\lvert(DX)^*\rvert=\sqrt{DX^2D^*}\).

\paragraph{R6. Acción y reloj.}
La sección de acción generada \(\hbar_{\rm ret}>0\) y el levantamiento
real de fase satisfacen
\begin{equation}
 S(\theta)=\hbar_{\rm ret}\theta.
 \label{g26:eq:accion-fase-es}
\end{equation}
Se conserva la velocidad constitutiva \(c>0\). Una vez obtenida la
longitud intensiva \(L\) de \eqref{g26:eq:regla-longitudinal-es}, el reloj
conjugado satisface \(\omega L=c\). Para \(\theta(t)=\omega t\),
la tasa de acción determina la escala energética
\begin{equation}
 E_{L}=\frac{dS}{dt}=\hbar_{\rm ret}\omega
   =\frac{\hbar_{\rm ret}c}{L}.
 \label{g26:eq:escala-energetica-es}
\end{equation}
El levantamiento y sus refinamientos permanecen asociados a la memoria
de la ruta.

\paragraph{R7. Lecturas del mismo carácter.}
Sobre \(Y=\mathsf U X\mathsf U^*>0\), se definen
\begin{equation}
 H_X=E_{L}Y,\qquad M_X=\frac{H_X}{c^2},\qquad
 \Lambda_X=LY^{-1}.
 \label{g26:eq:lecturas-conjugadas-es}
\end{equation}
La energía, la masa y la longitud conjugada conservan el mismo carácter
y el mismo transporte. La inversa actúa en el sector positivo; el sector
nulo conserva su tratamiento propio.

\paragraph{R8. Radio reducido y secciones dimensionales.}
La identificación física de esta realización es
\begin{equation}
 R_X=\frac{G}{c^2}M_X.
 \label{g26:eq:radio-reducido-es}
\end{equation}
Así, \(R_X\) representa el radio gravitatorio reducido. El radio de
Schwarzschild correspondiente es \(2R_X\). Se mantienen conjuntamente
\(L_*\), la sección de acción y la sección constitutiva de velocidad;
un cambio de unidades transforma sus coordenadas según sus dimensiones.

\paragraph{Procedencia de las restricciones.}
El transporte de Hadamard, la conservación del archivo, la inscripción,
la acción de fase y la reciprocidad poseen antecedentes en el corpus.
El registro ampliado \eqref{g26:eq:omega-rigidez-es} y las composiciones
\eqref{g26:eq:regla-longitudinal-es}--\eqref{g26:eq:metrica-radial-es} se
incorporan explícitamente en la construcción reunida. R4 y R5 son reglas
constitutivas de esta realización; la demostración siguiente establece
su consecuencia conjunta con los restantes antecedentes. R8 precisa
la identificación física que fija el significado del coeficiente.

\subsection{Teorema de determinación simultánea}

\paragraph{Teorema.}
Para los mismos datos generados \(B\), \((P_i)\), \(P\),
\(\alpha\), \(L_*\), \(\hbar_{\rm ret}\), \(c\) y el mismo carácter
\(X\), las restricciones R1--R8 determinan unívocamente la inscripción,
el transporte longitudinal y la respuesta radial. Para todos los
caracteres positivos admisibles, el coeficiente \(G\) es común. Se cumple
\begin{align}
 \mathcal E&=\Delta_\Omega,
 &\Delta_\Omega(T)&=\sum_{i\in\Omega}P_iTP_i,\label{g26:eq:delta-unica-es}\\
 r_N&=\frac{N-1}{4N}=\frac{5759}{23040},
 &L&=L_*\alpha^{16}r_N,\label{g26:eq:longitud-unica-es}\\
 D&=L\mathsf U,
 &R_X&=L\mathsf U X\mathsf U^*,\label{g26:eq:radial-unica-es}\\
 R_X\Lambda_X&=L^2I,
 &H_X\Lambda_X&=\hbar_{\rm ret}cI.\label{g26:eq:productos-radiales-es}
\end{align}
El acoplamiento queda determinado por
\begin{equation}
 \boxed{G=\frac{c^3L^2}{\hbar_{\rm ret}}
 =\frac{c^3L_*^2}{\hbar_{\rm ret}}
   \left(\frac{5759}{23040}\right)^2\alpha^{32}.}
 \label{g26:eq:g-rigido-es}
\end{equation}

\paragraph{Demostración: unicidad de la inscripción.}
Sean \(E_{ij}=e_ie_j^*\) las unidades matriciales. Por bimodularidad,
\[
 \mathcal E(E_{ij})=P_i\mathcal E(E_{ij})P_j.
\]
Para \(i\ne j\), la pertenencia de \(\mathcal E(E_{ij})\) al álgebra
diagonal implica \(P_i\mathcal E(E_{ij})P_j=0\). Para \(i=j\),
la fijación de \(P_i\) da \(\mathcal E(E_{ii})=P_i\).
La linealidad determina entonces \eqref{g26:eq:delta-unica-es} sobre todas
las matrices. La inscripción marcada queda así fijada por sus
propiedades algebraicas.

\paragraph{Demostración: coeficiente de retorno y archivo.}
La normalización de \(B\) y \(P^2=P=P^*\) proporcionan
\[
 C_0C_0^*=\tfrac14(I-P),\qquad M_0M_0^*=\tfrac14P.
\]
Como \(P_iPP_i=|u_i|^2P_i=P_i/N\), se obtiene
\begin{equation}
 \Delta_\Omega(C_0C_0^*)=\frac{N-1}{4N}I,
 \qquad
 \Delta_\Omega(M_0M_0^*)=\frac{1}{4N}I.
 \label{g26:eq:retorno-y-memoria-es}
\end{equation}
Ambas lecturas suman \(I/4\). La primera es escalar, de modo que
cualquier distribución normalizada \((p_i)\) sobre las incidencias,
con \(p_i\geq0\) y \(\sum_i p_i=1\), produce
\(\sum_i p_i r_N=r_N\). Se preserva, por tanto, el mismo coeficiente
bajo una ponderación diagonal arbitraria.

\paragraph{Demostración: longitud y métrica.}
La sustitución de \eqref{g26:eq:retorno-y-memoria-es} en R4 da
\(D=L\mathsf U\), con \(L\) como en \eqref{g26:eq:longitud-unica-es}.
La unitariedad de \(\mathsf U\) implica \(DD^*=L^2I\).
Por R5, las formas cuadráticas de \(R_X^2\) y \(DX^2D^*\) coinciden.
La identidad de polarización determina
\[
 R_X^2=DX^2D^*=L^2\mathsf U X^2\mathsf U^*.
\]
El operador \(L\mathsf U X\mathsf U^*\) es positivo y su cuadrado
es el operador del segundo miembro. La unicidad de la raíz cuadrada
positiva demuestra \eqref{g26:eq:radial-unica-es}. En particular,
\(R_I=LI\).

\paragraph{Demostración: acción, reciprocidad y acoplamiento.}
R6 y R7 dan
\[
 H_X=\frac{\hbar_{\rm ret}c}{L}Y,
 \qquad M_X=\frac{\hbar_{\rm ret}}{cL}Y,
 \qquad \Lambda_X=LY^{-1}.
\]
La multiplicación produce ambos productos de
\eqref{g26:eq:productos-radiales-es}. A continuación, R8 equivale a
\[
 LY=\frac{G\hbar_{\rm ret}}{c^3L}Y.
\]
La invertibilidad de \(Y\) determina \eqref{g26:eq:g-rigido-es}.
La misma sustitución verifica la igualdad para cada carácter positivo.
Si \(G_1\) y \(G_2\) satisfacen la identidad con los mismos datos,
\((G_1-G_2)Y=0\), y por tanto \(G_1=G_2\).
Esta conclusión distingue la respuesta dependiente del estado
\(R_X\) del coeficiente común \(G\). \hfill\(\square\)

\subsection{Reciprocidad y rigidez de la fase}

El antecedente potencia--logarítmico utiliza, para \(x>0\),
\[
 \ell_p(x)=\frac{x^p-1}{p}\quad(p\ne0),
 \qquad \ell_0(x)=\log x.
\]
La sustitución directa demuestra
\(\ell_{-p}(x^{-1})=-\ell_p(x)\). En el dominio positivo de sus
inversas, esta identidad equivale a
\[
 \exp_p(s)\exp_{-p}(-s)=1.
\]
La composición espectral sobre un carácter positivo transmite este
producto al par \(LY\), \(LY^{-1}\). En la realización actual, R5
puede expresarse equivalentemente, una vez fijadas R1--R4 y R7, por
\(R_X\Lambda_X=L^2I\): multiplicar por \(\Lambda_X^{-1}\) determina
\(R_X=L\mathsf U X\mathsf U^*\), y esta expresión satisface R5.

Los dos productos de \eqref{g26:eq:productos-radiales-es} dan además una
forma directa de la rigidez:
\begin{equation}
 R_X=\frac{L^2}{\hbar_{\rm ret}c}H_X.
 \label{g26:eq:proporcion-operadores-es}
\end{equation}
Para una reescala positiva \(s\), el par
\((sR_X,\Lambda_X/s)\) conserva el producto radial; el producto
\(H_X(\Lambda_X/s)=\hbar_{\rm ret}cI/s\), con energía y acción
fijadas, exige \(s=1\). La conservación conjunta determina la escala.
Asimismo, cualquier operador radial positivo con la misma métrica R5
coincide con \(R_X\) por la unicidad de la raíz positiva.

La acción levantada también posee una determinación exacta. Para
\(\theta\ne0\), la igualdad \(S/\hbar_{\rm ret}=\theta\) fija \(S\).
Incluso cuando se conservan las rotaciones de todas las subdivisiones
\(\theta/9^n\), una reescala real fija \(a\) que satisfaga
\[
 \exp(-ia\theta/9^n)=\exp(-i\theta/9^n)
 \qquad(n\geq0)
\]
verifica \((a-1)\theta/9^n\in2\pi\mathbb Z\). Esta sucesión tiende
a cero; para \(n\) suficientemente grande su único valor posible es
cero. De \(\theta\ne0\) se deduce \(a=1\).

\subsection{Escalas de Planck y especialización del carácter unitario}

Escribimos \(L_\alpha=L\); en la sección de retorno, \(E_{P,\rm ret}=E_L\).

La longitud construida $L_\alpha>0$, la sección de acción $\hbar_a>0$ y la
velocidad $c>0$ determinan las escalas

\[
t_P=\frac{L_\alpha}{c},\qquad
m_{P,a}=\frac{\hbar_a}{cL_\alpha},\qquad
E_{P,a}=\frac{\hbar_a c}{L_\alpha}.
\]

En efecto, sustituir $G_a=c^3L_\alpha^2/\hbar_a$ en las expresiones
$\sqrt{\hbar_aG_a/c^5}$, $\sqrt{\hbar_ac/G_a}$ y
$\sqrt{\hbar_ac^5/G_a}$ da, respectivamente, esas tres identidades.
La positividad fija en cada caso la raíz. La longitud antecede al
reconocimiento de estas escalas; sus expresiones convencionales son
identidades de la realización ya construida.

Para un modo de carácter $y>0$, las lecturas másica y longitudinal son

\[
m(y)=m_{P,a}y,\qquad r_g(y)=L_\alpha y,\qquad
\bar\lambda(y)=\frac{L_\alpha}{y}.
\]

Por tanto, $r_g(y)\bar\lambda(y)=L_\alpha^2$ para todos los modos
positivos. La igualdad de ambas longitudes caracteriza el modo $y=1$,
puesto que $y=y^{-1}$ y $y>0$ implican $y=1$. En este modo, la masa es
$m_{P,a}$ y ambas longitudes coinciden con $L_\alpha$. Así se distingue
la reciprocidad general de su especialización en el carácter unitario.

El paso cronológico y el tiempo de Planck satisfacen
$t_0=(\pi_{\rm HMT}/54)t_P$. El período completo de fase de 108 pasos es
$108t_0=2\pi_{\rm HMT}t_P$ y su frecuencia angular es
$\omega_{108}=1/t_P=c/L_\alpha$. Un horizonte de radio
$R_h=\zeta_hL_\alpha$, con $\zeta_h>0$, conserva su variable geométrica
propia. El radio de Schwarzschild asociado a una masa es $2r_g$.

Si se comparan las dos secciones de acción con $L_\alpha$ y $c$ fijados,
el cociente $\rho=\hbar_b/\hbar_a$ da
$G_b=G_a/\rho$, $m_{P,b}=\rho m_{P,a}$ y
$E_{P,b}=\rho E_{P,a}$; $L_\alpha$ y $t_P$ permanecen iguales.
Esta comparación entre secciones conserva la distinción respecto de una
evolución temporal de las constantes.

\subsection{Covarianza del archivo y transformación de unidades}

Sean \(T:V_{\rm in}\to V'_{\rm in}\) y
\(S:V_{\rm out}\to V'_{\rm out}\) unitarios. El transporte conjunto
de los datos es
\[
 B'=SBT^*,\quad P'=SPS^*,\quad P_i'=SP_iS^*,\quad X'=TXT^*.
\]
La expectativa transportada
\(\mathcal E'(A')=S\mathcal E(S^*A'S)S^*\) fija \(P_i'\)
y es bimodular en el álgebra marcada transformada. Por sustitución,
\begin{align*}
 \mathsf U'&=S\mathsf UT^*,&D'&=SDT^*,\\
 R'_{X'}&=SR_XS^*,&H'_{X'}&=SH_XS^*,&
 \Lambda'_{X'}&=S\Lambda_XS^*.
\end{align*}
La conjugación conserva las igualdades de R1--R8 y el mismo coeficiente
\(G\). Las marcas, orientaciones y registros se transportan con las
bases. Las permutaciones de incidencias son un caso particular.

La incorporación de un espacio auxiliar de memoria \(\mathcal F\)
transporta \(P\) a \(P\otimes I_{\mathcal F}\),
\(\mathsf U\) a \(\mathsf U\otimes I_{\mathcal F}\) y
\(\Delta_\Omega\) a \(\Delta_\Omega\otimes\operatorname{id}\).
La inscripción de \(C_0C_0^*\otimes I_{\mathcal F}\) es
\(r_N I\otimes I_{\mathcal F}\); el coeficiente conserva el cardinal
marcado original.

Para precisar la covarianza dimensional, sean \(\ell\), \(\tau\) y
\(m\) los factores positivos de las nuevas unidades de longitud, tiempo
y masa respecto de las anteriores. Las coordenadas numéricas satisfacen
\[
 L'=L/\ell,\qquad c'=c\tau/\ell,\qquad
 \hbar'_{\rm ret}=\hbar_{\rm ret}\tau/(m\ell^2).
\]
Así,
\[
 \frac{(c')^3(L')^2}{\hbar'_{\rm ret}}
 =G\frac{m\tau^2}{\ell^3},
\]
que es la transformación de una magnitud de dimensión
\(\mathrm{longitud}^3\mathrm{masa}^{-1}\mathrm{tiempo}^{-2}\).
La sección \(L_*\) se transforma como \(L\), mientras \(\alpha\)
y \(r_N\) permanecen adimensionales.

\subsection{Eliminación exacta de memoria acoplada}

Sean \(\mathcal K\) y \(\mathcal H\) espacios de Hilbert finitos,
\(W=W^*>0\) un peso global y
\(\mathcal C:\mathcal K\to\mathcal H\) sobreyectiva. La acción
de los residuos \(r\), con defecto de frontera \(z\), es
\[
 \mathcal S(r)=\mathfrak a\,r^*Wr,
 \qquad \mathcal Cr=z,\qquad \mathfrak a>0.
\]
El peso global conserva sus bloques cruzados entre incidencias. Defínase
\(\mathcal R=\mathcal CW^{-1}\mathcal C^*\). Para \(v\ne0\),
la inyectividad de \(\mathcal C^*\) da
\[
 v^*\mathcal Rv=\|W^{-1/2}\mathcal C^*v\|^2>0.
\]
Por tanto, \(\mathcal R\) es invertible. El residuo de mínima acción es
\[
 r_*(z)=W^{-1}\mathcal C^*\mathcal R^{-1}z.
\]
Se verifica \(\mathcal Cr_*=z\). Cada residuo admisible se escribe
\(r=r_*+k\), con \(\mathcal Ck=0\), y
\(k^*Wr_*=k^*\mathcal C^*\mathcal R^{-1}z=0\). En consecuencia,
\begin{equation}
 \mathcal S(r)=\mathfrak a\,z^*\mathcal R^{-1}z
             +\mathfrak a\,k^*Wk.
 \label{g26:eq:accion-residuo-es}
\end{equation}
La positividad de \(W\) prueba la unicidad del mínimo. El par
\((z,k)\) conserva y reconstruye el residuo completo.

Para una ruta con transportes \(U_j\), los residuos
\(r_j=f_j-U_jf_{j-1}\) se reconstruyen mediante
\(f_j=U_jf_{j-1}+r_j\). La aplicación \(\mathcal C\) reúne los
transportes posteriores de cada residuo hacia el extremo. Si sus
bloques son \(\mathcal C_j\), entonces
\[
 \mathcal R=\sum_{j,k}\mathcal C_j(W^{-1})_{jk}\mathcal C_k^*.
\]
Esta expresión conserva las correlaciones entre tramos de la ruta.
Para un transporte longitudinal invertible \(D\), la frontera
\(y=Dz\) tiene respuesta
\[
 \mathcal R_L=D\mathcal RD^*,\qquad
 \mathcal S_{{\rm ef},L}(y)
 =\mathfrak a\,y^*(D\mathcal RD^*)^{-1}y
 =\mathfrak a\,z^*\mathcal R^{-1}z.
\]
Se obtiene la misma acción al eliminar primero la memoria o transportar
primero la frontera. El prefactor conserva la normalización de la acción
fuente; cuando esa acción lleva \(\mathfrak a=\hbar_{\rm ret}\), la
reducción transmite la misma sección.

\subsection{Reducción conjunta de radio, energía y longitud conjugada}

Escríbase el mismo carácter positivo en una descomposición de frontera
y memoria:
\[
 Y=\begin{pmatrix}A&F\\F^*&C\end{pmatrix},\qquad C>0,
 \qquad Y_{\rm ef}=A-FC^{-1}F^*.
\]
La identidad
\[
 \begin{pmatrix}b\\y\end{pmatrix}^{\!*}Y
 \begin{pmatrix}b\\y\end{pmatrix}
 =b^*Y_{\rm ef}b+\eta^*C\eta,
 \qquad \eta=y+C^{-1}F^*b,
\]
prueba \(Y_{\rm ef}>0\), caracteriza la extensión de mínima energía y
conserva la reconstrucción \(y=\eta-C^{-1}F^*b\). Para \(s>0\),
la sustitución en el complemento de Schur da
\(\operatorname{Schur}(sY)=sY_{\rm ef}\). En consecuencia,
\[
 R_{\rm ef}=LY_{\rm ef},\qquad
 H_{\rm ef}=E_{L}Y_{\rm ef}.
\]
Para calcular la longitud conjugada comprimida, se resuelve
\(Y(x,y)=(f,0)\). La segunda ecuación da
\(y=-C^{-1}F^*x\), y la primera da
\(x=Y_{\rm ef}^{-1}f\). El bloque de frontera de \(Y^{-1}\) es,
por tanto, \(Y_{\rm ef}^{-1}\). Así,
\begin{equation}
 \Lambda_{\rm b}=LY_{\rm ef}^{-1},\qquad
 R_{\rm ef}\Lambda_{\rm b}=L^2I,\qquad
 R_{\rm ef}=\frac{L}{E_{L}}H_{\rm ef}.
 \label{g26:eq:schur-reciproco-es}
\end{equation}
La reducción actúa sobre los mismos bloques en las dos lecturas y
conserva el residuo \(\eta\). El carácter efectivo puede depender de
la memoria de forma matricial y el coeficiente \(L/E_{L}\) permanece fijo.

\subsection{Memoria dinámica y norma temporal inducida}

La eliminación dinámica conserva una dependencia espectral adicional.
En una representación estacionaria, sea
\[
 H=\begin{pmatrix}H_{\rm bb}&V\\V^*&H_{\rm ii}\end{pmatrix},
 \qquad H_{\rm ii}>0.
\]
Para \(z\) en el dominio resolvente pertinente, la eliminación de la
componente interior en \((H-zI)(b,y)=(f,0)\) da
\begin{equation}
 \mathcal K_H(z)=H_{\rm bb}-zI
       -V(H_{\rm ii}-zI)^{-1}V^*.
 \label{g26:eq:resolvente-memoria-es}
\end{equation}
Cuando existen ambos inversos, el bloque de frontera de
\((H-zI)^{-1}\) es \(\mathcal K_H(z)^{-1}\). Para
\(|z|<\lambda_{\min}(H_{\rm ii})\), la serie convergente del
resolvente proporciona
\begin{align}
 \mathcal K_H(z)&=H_{\rm est}-zZ-z^2VH_{\rm ii}^{-3}V^*-\cdots,
 \label{g26:eq:expansion-resolvente-es}\\
 H_{\rm est}&=H_{\rm bb}-VH_{\rm ii}^{-1}V^*,\qquad
 Z=I+VH_{\rm ii}^{-2}V^*\geq I.
 \label{g26:eq:norma-z-es}
\end{align}
La positividad se sigue de
\(Z-I=(H_{\rm ii}^{-1}V^*)^*(H_{\rm ii}^{-1}V^*)\).
La extensión estática \((b,-H_{\rm ii}^{-1}V^*b)\) tiene norma
cuadrada \(b^*Zb\); la energía y la norma temporal proceden de la
misma extensión. El parámetro físico es \(z=\hbar_{\rm ret}\omega\).

Al primer orden temporal, \(T=Z^{1/2}\) y \(\psi=Tb\) transforman
la pareja \((H_{\rm est},Z)\) en
\((T^{-*}H_{\rm est}T^{-1},I)\). La sección
\(\hbar_{\rm ret}\) permanece fijada. Si \(T\) depende del tiempo,
la transformación conserva también el término \(\dot T b\) de
\(\dot\psi=\dot T b+T\dot b\).

La proporcionalidad radial se conserva en el resolvente completo.
Escribiendo \(\chi_{\rm rad}=L/E_{L}\) y \(R=\chi_{\rm rad} H\), la sustitución de
cada bloque demuestra
\begin{equation}
 \mathcal K_R(\chi_{\rm rad} z)=\chi_{\rm rad}\mathcal K_H(z).
 \label{g26:eq:resolvente-proporcional-es}
\end{equation}
Esta identidad transporta simultáneamente el operador y el parámetro
espectral, y conserva cada orden de la memoria dinámica. Para las formas
efectivas de \eqref{g26:eq:schur-reciproco-es}, la normalización canónica
requiere el transporte dual de la longitud inversa:
\begin{align*}
 H_{\rm can}&=T^{-*}H_{\rm ef}T^{-1},&
 R_{\rm can}&=T^{-*}R_{\rm ef}T^{-1},&
 \Lambda_{\rm can}&=T\Lambda_{\rm b}T^*.
\end{align*}
Multiplicar estas expresiones da
\[
 R_{\rm can}=\frac{L}{E_{L}}H_{\rm can},\qquad
 R_{\rm can}\Lambda_{\rm can}=L^2I.
\]
La demostración admite \(T\), \(R_{\rm ef}\) y
\(\Lambda_{\rm b}\) no conmutativos. La norma temporal inducida y
la longitud dual conservan conjuntamente el mismo \(G\).

\subsection{Refinamiento, dominios espectrales y alcance}

Sean \(J\) y \(K\) inscripciones isométricas entre dos niveles que
satisfagan \(X'J=JX\) y \(\mathsf U'J=K\mathsf U\),
con la misma sección intensiva \(L\). Entonces
\[
 R'K=L\mathsf U'X'\mathsf U'^*K
     =KL\mathsf U X\mathsf U^*=KR,
\]
y la misma cuenta demuestra \(H'K=KH\). El entrelazamiento de las
inversas se formula en los dominios correspondientes. Cuando el sistema
compatible posee un límite autoadjunto positivo e inyectivo \(X\),
el cálculo espectral define \(Y=\mathsf U X\mathsf U^*\),
\(R=LY\), \(H=E_{L}Y\) y \(\Lambda=LY^{-1}\). La proporcionalidad
\(R=(L/E_{L})H\) vale en \(\operatorname{Dom}(Y)\), mientras
\(R\Lambda=L^2I\) y \(H\Lambda=\hbar_{\rm ret}cI\) valen en
\(\operatorname{Dom}(Y^{-1})\), con prolongación acotada de sus
productos. Los cortes espectrales
\(\mathbf1_{[\varepsilon,M]}(Y)\) verifican las identidades antes
de pasar al límite en estos dominios. Una longitud de arista refinada
\(L/9\) incorpora su factor de escala propio; la sección intensiva
se conserva según el entrelazamiento declarado.

La clasificación distingue las transformaciones que preservan la
realización de las que modifican una regla constitutiva. Las ponderaciones
diagonales, las equivalencias unitarias, el archivo auxiliar y la reducción
conjunta conservan \(G\). El cambio de resolución marcada modifica R1;
la sustitución de \(r_N\) por \(\sqrt{r_N}\), de
\(\alpha^{16}\) por su cuadrado o de la componente inscrita por la
norma total modifica R4. La reescala aislada del radio cambia R5,
y la reescala aislada de la acción cambia R6. R8 mantiene la distinción
entre radio reducido, radio de Schwarzschild y cambio de unidades.

Se obtiene así la unicidad matemática de la realización R1--R8: los datos
generados y el carácter fijan la respuesta, y la composición común fija
\(G\). Las reglas longitudinal y métrica figuran explícitamente en
esta conclusión como composiciones incorporadas; la identificación
gravitatoria conserva su estatuto físico. La prueba abarca la clase
declarada con sus definiciones constitutivas y sus transformaciones
compatibles. Las comprobaciones racionales del desarrollo verifican
casos finitos de las identidades expuestas; la demostración general
reside en las ecuaciones y argumentos anteriores.
